\documentclass{article}
\usepackage{amsmath,amssymb}
\usepackage[vocab={numbers,arithmetic,congruences}]{leanify}
\newtheorem{lemma}{Lemma}
\begin{document}

\section{Notation}

We write $\varphi$ for Euler's totient function.
We write $\left(\frac{a}{p}\right)$ for the Legendre symbol of $a$ modulo $p$.
We write $L(a,p)$ for the Legendre symbol of $a$ modulo $p$.
We write $\operatorname{inv}_n(a)$ for the inverse of $a$ modulo $n$.
We write $g(n)$ for the least primitive root modulo $n$.
We write $N(n)$ for the number of quadratic residues modulo $n$.

\section{Congruences}

\begin{lemma}\label{euler}
Let $n$ be a positive integer and let $a$ be an integer coprime to $n$. Then $a^{\varphi(n)}$ is congruent to $1$ modulo $n$.
\end{lemma}

\begin{lemma}\label{fermat}
Let $p$ be a prime and let $a$ be an integer. Then $a^p$ is congruent to $a$ modulo $p$.
\end{lemma}

\begin{lemma}\label{residues}
$2$ is a quadratic residue modulo $7$, and $3$ is a quadratic nonresidue modulo $7$. Moreover, $L(2,7)=1$ and $L(3,7)=-1$.
\end{lemma}

\begin{lemma}\label{count}
For every prime $p>2$, $N(p)=\frac{p-1}{2}$.
\end{lemma}

\begin{lemma}\label{roots}
$3$ is a primitive root modulo $7$, and $g(7)=3$ and $\operatorname{inv}_7(3)=5$.
\end{lemma}

\begin{lemma}\label{criterion}
Let $p>2$ be a prime and let $a$ be an integer coprime to $p$. Then $a^{\frac{p-1}{2}}$ is congruent to $L(a,p)$ modulo $p$.
\end{lemma}

\begin{lemma}\label{fermat-little}
Let $p$ be a prime and let $a$ be an integer not divisible by $p$. Then $a^{p-1}\equiv1\pmod p$.
\end{lemma}

\begin{lemma}\label{prime-divisor}
Let $n\ge2$ be an integer. Then $n$ has a prime divisor.
\end{lemma}

\begin{lemma}\label{bezout}
Let $a$ and $b$ be coprime positive integers. Then there are integers $x$ and $y$ with $ax+by=1$.
\end{lemma}

\begin{lemma}\label{six}
For every integer $n\ge0$, the number $n^3-n$ is divisible by $6$.
\end{lemma}

\end{document}
